\chapter{Subsucesión monótona}

\section{Introducción: El teorema de la subsucesión monótona}
%-----------------------------------------------------------------------

\begin{claim}
Toda sucesión infinita de números reales (o racionales) tiene una subsucesión
\emph{monótona} (no decreciente) o \emph{antítona} (no creciente).
\end{claim}

Este resultado fundamental es una piedra angular del análisis real y proporciona
una visión profunda de la estructura de las sucesiones. En este capítulo
demostraremos la mitad del teorema: que las sucesiones sin picos no acotados
poseen una subsucesión monótona.

La estrategia gira en torno a una hermosa dicotomía. Definiremos qué significa
que una sucesión tenga \emph{picos no acotados}: intuitivamente, posiciones desde
las cuales la sucesión nunca vuelve a subir, y de las que existen arbitrariamente
muchas. La idea clave es que si una sucesión no tiene esta propiedad, podemos
construir una subsucesión creciente evitando hábilmente los picos.

\section{Nuevas definiciones}

\subsection{Picos en una sucesión}

Consideremos una sucesión $a : \mathbb{N} \to X$, donde $X$ puede ser
$\mathbb{Q}$ o $\mathbb{R}$ (o cualquier conjunto linealmente ordenado).
Decimos que la posición $n$ es un \emph{pico} si desde ese punto podemos mirar
hacia adelante y ver que todos los valores futuros son menores o iguales que
$a(n)$.

\begin{definicion}
Para una sucesión $a : \mathbb{N} \to X$ y un índice $n : \mathbb{N}$, decimos
que $n$ es un \textbf{pico} de $a$ si:
\[
  \mathrm{IsAPeak}(a,\, n) \;:=\; \forall\, m > n,\quad a(m) \leq a(n).
\]
En otras palabras, parados en la posición $n$, podemos mirar hacia abajo
todos los valores futuros.
\end{definicion}

\subsection{Picos no acotados}

Una sucesión puede tener muchos picos, pocos, o incluso ninguno. Nos interesan
especialmente las sucesiones cuyos picos son \emph{no acotados}: sin importar
qué tan lejos miremos, siempre hay otro pico aún más adelante.

\begin{definicion}
Una sucesión $a : \mathbb{N} \to X$ tiene \textbf{picos no acotados} si:
\[
  \mathrm{UnBddPeaks}(a) \;:=\; \forall\, k,\; \exists\, n > k,\;
  \mathrm{IsAPeak}(a,\, n).
\]
Esto afirma que para cualquier cota $k$, existe una posición de pico $n$ más
allá de $k$. El conjunto de picos no tiene cota superior.
\end{definicion}

\subsection{Expresiones condicionales: \texttt{if-then-else}}

En Lean, podemos definir valores de forma condicional. La sintaxis es:

\begin{minted}{lean4}
if h : condicion then valor_si_verdadero else valor_si_falso
\end{minted}

El \texttt{h :} opcional enlaza una prueba de la condición (cuando es
verdadera) o su negación (cuando es falsa), disponible en la rama respectiva.
Es una característica poderosa que usaremos para construir nuestra subsucesión.

\section{Nuevo teorema: \texttt{lt\_of\_not\_ge}}

Una de las propiedades fundamentales de los conjuntos ordenados es que las
comparaciones son decidibles: para cualesquiera dos elementos $m, n$, se tiene
$m \leq n$ o bien $n < m$.

\begin{theorem}[\texttt{lt\_of\_not\_ge}]
Si $\lnot(m \leq n)$, entonces $n < m$.
\end{theorem}

Este hecho, aparentemente simple, es crucial para nuestra construcción: nos
permite extraer desigualdad estricta a partir del fracaso de la desigualdad
débil.

\section{Teorema principal: Subsucesión monótona a partir de picos acotados}

Llegamos ahora a nuestro resultado principal. El teorema establece que si una
sucesión no tiene picos no acotados, entonces debe poseer una subsucesión
monótona creciente.

\begin{theorem}[\texttt{MonotoneSubseq\_of\_BddPeaks}]
Sea $a : \mathbb{N} \to X$ una sucesión (donde $X$ es $\mathbb{Q}$ o
$\mathbb{R}$). Si $a$ no tiene picos no acotados, entonces existe una
subsucesión $\sigma : \mathbb{N} \to \mathbb{N}$ tal que:
\begin{itemize}
  \item $\sigma$ es estrictamente creciente (\texttt{a Subseq}).
  \item $a \circ \sigma$ es monótona (no decreciente).
\end{itemize}
\end{theorem}

\subsection{Visión estratégica de la demostración}

La demostración se desarrolla en varias etapas.

\medskip
\textbf{Etapa 1: Desempacar la hipótesis.}
La negación de $\mathrm{UnBddPeaks}(a)$ significa que existe alguna cota $k$
tal que no hay picos más allá de $k$. Usando \texttt{choose}, extraemos:
\begin{itemize}
  \item Una cota $k : \mathbb{N}$.
  \item Para cada $n > k$, un testigo $\tau(n)$ con $\tau(n) > n$ y
        $\lnot\mathrm{IsAPeak}(a, \tau(n))$.
\end{itemize}

\medskip
\textbf{Etapa 2: Construir una sucesión auxiliar.}
Definimos $\tau' : \mathbb{N} \to \mathbb{N}$ que crece más rápido que la
identidad:
\[
  \tau'(n) :=
  \begin{cases}
    n + 1   & \text{si } n \leq k, \\
    \tau(n) & \text{si } n > k.
  \end{cases}
\]
El primer caso cubre índices hasta la cota $k$; el segundo usa los testigos
de no-pico para índices más allá de $k$.

\medskip
\textbf{Etapa 3: Construir la órbita.}
Usando la construcción de órbita del capítulo anterior, definimos:
\[
  \sigma(n) := (\tau')^{[n]}(k+1).
\]
Como $\tau'$ crece más rápido que la identidad, el teorema
\texttt{Subseq\_of\_Iterate} garantiza que $\sigma$ es estrictamente creciente.

\medskip
\textbf{Etapa 4: Demostrar la monotonía.}
La observación clave es que $\sigma(n) > k$ para todo $n$, de modo que
$\sigma(n)$ nunca es un pico. Como $\sigma(n)$ no es un pico, existe
$m > \sigma(n)$ con $a(m) > a(\sigma(n))$. Por nuestra construcción,
$\sigma(n+1) = \tau'(\sigma(n)) > \sigma(n)$ es precisamente ese $m$, lo que
da $a(\sigma(n)) \leq a(\sigma(n+1))$.

\subsection{Demostración en Lean}

\begin{minted}{lean4}
Statement MonotoneSubseq_of_BddPeaks (a : N -> X)
    (ha : ¬ UnBddPeaks a)
    : ∃ σ, Subseq σ ∧ Monotone (a ∘ σ) := by
  change ¬ (∀ k, ∃ n > k, ∀ m > n, a m ≤ a n) at ha
  push_neg at ha
  choose k hk using ha
  choose τ τ_gt hτ using hk
  let τ' : N → N := fun n =>
    if h : n ≤ k then n + 1 else τ n (lt_of_not_ge h)
  have τ'_eq : ∀ n, τ' n =
      if h : n ≤ k then n + 1 else τ n (lt_of_not_ge h) := by
    intro n; rfl
  have τ'_gt : ∀ n, n < τ' n := by
    intro n
    by_cases hn : n ≤ k
    · rewrite [τ'_eq]; bound
    · rewrite [τ'_eq]; bound
  let σ : N → N := fun n ↦ τ'^[n] (k + 1)
  have σ_eq : ∀ n, σ n = τ'^[n] (k + 1) := by intro n; rfl
  have hσ : ∀ n, k < σ n := by
    intro n
    induction' n with n hn
    · rewrite [σ_eq]; bound
    · rewrite [σ_eq]
      rewrite [← show τ' (τ'^[n] (k + 1)) = τ'^[n + 1] (k + 1)
        by apply succ_iterate]
      rewrite [← σ_eq]
      specialize τ'_gt (σ n)
      linarith [τ'_gt, hn]
  use σ
  split_ands
  · apply Subseq_of_Iterate
    apply τ'_gt
  · apply Monotone_of_succ
    intro n
    specialize hσ n
    specialize hτ (σ n) hσ
    change a (σ n) ≤ a (σ (n + 1))
    rewrite [show σ (n + 1) = τ'^[n + 1] (k + 1) by apply σ_eq]
    rewrite [← show τ' (τ'^[n] (k + 1)) = τ'^[n + 1] (k + 1)
      by apply succ_iterate]
    rewrite [← show σ n = τ'^[n] (k + 1) by apply σ_eq]
    rewrite [τ'_eq]
    bound
\end{minted}

\subsection{Demostración en lenguaje natural}

\begin{proof}
Supongamos que $a : \mathbb{N} \to X$ no tiene picos no acotados.
Construiremos una subsucesión estrictamente creciente a lo largo de la cual
$a$ es monótona.

\medskip
\textbf{Paso 1: Extraer la cota.}
Como $\lnot\mathrm{UnBddPeaks}(a)$, la negación de
$\forall k,\,\exists n>k,\,\mathrm{IsAPeak}(a,n)$ nos proporciona:
\[
  \exists k,\;\forall n > k,\;\lnot\mathrm{IsAPeak}(a,n).
\]
Por la táctica \texttt{choose}, obtenemos un $k:\mathbb{N}$ específico.

\medskip
\textbf{Paso 2: Extraer los testigos.}
Para cada $n > k$, como $\lnot\mathrm{IsAPeak}(a,n)$:
\[
  \lnot(\forall m>n,\;a(m)\leq a(n))
  \;\equiv\;
  \exists m>n,\;a(m)>a(n).
\]
Usando \texttt{choose} nuevamente, obtenemos una función $\tau$ tal que
$\tau(n)>n$ y $a(\tau(n))>a(n)$ para todo $n>k$.

\medskip
\textbf{Paso 3: Definir la sucesión auxiliar.}
Definimos $\tau':\mathbb{N}\to\mathbb{N}$ por:
\[
  \tau'(n):=\begin{cases}n+1&\text{si }n\leq k,\\\tau(n)&\text{si }n>k,\end{cases}
\]
usando \texttt{lt\_of\_not\_ge} para convertir $\lnot(n\leq k)$ en $k<n$.
Se verifica que $n < \tau'(n)$ en ambos casos.

\medskip
\textbf{Paso 4: Construir la órbita.}
Definimos $\sigma:\mathbb{N}\to\mathbb{N}$ por
$\sigma(n):=(\tau')^{[n]}(k+1)$.
Por \texttt{Subseq\_of\_Iterate}, $\sigma$ es estrictamente creciente.

\medskip
\textbf{Paso 5: Verificar la desigualdad clave.}
Por inducción, $\sigma(n)>k$ para todo $n$:
\begin{itemize}
  \item \textit{Caso base}: $\sigma(0)=k+1>k$.
  \item \textit{Paso inductivo}: si $\sigma(n)>k$, entonces
        $\sigma(n+1)=\tau'(\sigma(n))>\sigma(n)>k$.
\end{itemize}

\medskip
\textbf{Paso 6: Demostrar la monotonía.}
Basta mostrar $a(\sigma(n))\leq a(\sigma(n+1))$ para todo $n$
(por \texttt{Monotone\_of\_succ}). Como $\sigma(n)>k$:
\begin{align*}
  &\sigma(n+1)=\tau'(\sigma(n))=\tau(\sigma(n)), \quad \text{y}\quad\\
  &a(\sigma(n))\leq a(\tau(\sigma(n)))=a(\sigma(n+1)).
\end{align*}

\textbf{Conclusión:} $\sigma$ es estrictamente creciente y $a\circ\sigma$ es
monótona. \qed
\end{proof}

\section{El panorama general}

Este teorema revela una dicotomía fundamental en el comportamiento de las
sucesiones. Toda sucesión cae en uno de dos grupos:

\begin{enumerate}
  \item \textbf{Sucesiones con picos no acotados:} tienen una subsucesión
        no creciente (antítona). (Contenido del ejercicio de tarea.)
  \item \textbf{Sucesiones con picos acotados:} tienen una subsucesión no
        decreciente (monótona). (Lo que acabamos de demostrar.)
\end{enumerate}

Juntos, estos resultados implicarán el \textbf{teorema de Bolzano--Weierstrass}
para sucesiones: toda sucesión infinita tiene una subsucesión de Cauchy.
Este resultado es fundamental porque:
\begin{itemize}
  \item Combinado con la acotación, garantiza subsucesiones convergentes.
  \item Provee una herramienta para extraer comportamiento regular de
        sucesiones aparentemente caóticas.
  \item Se conecta con la compacidad en espacios métricos.
  \item Subyace a muchas demostraciones de existencia en análisis.
\end{itemize}

\subsection{El papel de la construcción \texttt{if-then-else}}

La definición condicional de $\tau'$ es esencial. Necesitábamos una sola
función que:
\begin{itemize}
  \item Crezca más rápido que la identidad en todas partes.
  \item Coincida con la función testigo $\tau$ para índices más allá de $k$.
  \item Esté bien definida para todos los índices, incluidos los $\leq k$
        donde $\tau$ no existe.
\end{itemize}

La construcción \texttt{if-then-else} maneja ambos casos de forma elegante,
permitiéndonos ensamblar una función globalmente definida a partir de un
testigo parcialmente definido. Esta técnica de «parchear» funciones es común
en el análisis constructivo.

\subsection{Por qué funcionan las órbitas}

La construcción $\sigma(n)=(\tau')^{[n]}(k+1)$ es perfecta aquí porque:
\begin{enumerate}
  \item Proporciona automáticamente una sucesión estrictamente creciente
        (por \texttt{Subseq\_of\_Iterate}).
  \item Nos mantiene en la región donde $\tau'$ coincide con $\tau$
        (más allá de $k$).
  \item Cada iteración nos lleva a un punto con valor estrictamente mayor.
  \item El punto de partida $k+1$ nos sitúa de inmediato en la región activa.
\end{enumerate}

Esto demuestra el poder del pensamiento en sistemas dinámicos aplicado al
análisis: viendo $\tau'$ como un sistema dinámico y estudiando sus órbitas,
extraemos el comportamiento monótono deseado.

\section{Perspectivas futuras}

En la tarea demostrarán el resultado complementario: las sucesiones con picos
no acotados tienen subsucesiones antítonas (no crecientes). La estrategia es
similar, pero involucra elegir los propios picos para formar la subsucesión.

Junto con los resultados sobre sucesiones acotadas, estos teoremas forman la
base para entender cuándo las sucesiones tienen subsucesiones convergentes,
pregunta central en el análisis real.

Las técnicas aquí desarrolladas ---usar \texttt{choose} para extraer testigos
de enunciados de existencia, definiciones condicionales para construir
funciones, e iteraciones de órbitas para construir subsucesiones--- son
herramientas fundamentales que utilizarán a lo largo del análisis.

%%% Local Variables:
%%% mode: LaTeX
%%% TeX-master: "Analisis_en_Lean"
%%% End:
